
\section{Implementación del prototipo}
\label{sec:implementacion}

Si la sección anterior describía el esqueleto, esta lo recubre de músculo.
Recorremos las piezas en el mismo orden en que intervienen en el flujo:
primero el detector, que decide qué se puede combinar; después el
transformador, que efectúa la combinación; luego el cálculo de la
complejidad cognitiva, con el que medimos el efecto; y, por último, el
pipeline experimental y el subsistema de evaluación de LLMs. 

\subsection{Detector}

La clase \texttt{NestedIfDetector} recorre el AST del método y visita todos
los \lstinline|if|, incluidos los anidados a cualquier profundidad. Para
cada \lstinline|if| candidato aplica en orden las precondiciones:

\begin{enumerate}
    \item \textbf{P1.} El \lstinline|if| externo no tiene rama \lstinline|else|.

    \item \textbf{P2--P3.} Su rama \lstinline|then| consiste en una única
    sentencia y esa sentencia es, a su vez, un \lstinline|if|. Se aceptan
    tanto el bloque con llaves de una sola sentencia como la forma sin
    llaves; en cualquier otro caso, el candidato se descarta por P2 o P3.

    \item \textbf{P4.} El \lstinline|if| interno tampoco tiene rama
    \lstinline|else|.

    \item \textbf{P5.} Ninguna de las dos condiciones contiene efectos
    colaterales: las cinco categorías descritas en \autoref{sec:precondiciones}.
\end{enumerate}

El detector ofrece dos modos de salida: uno devuelve solo las oportunidades
aceptadas y otro (\lstinline|detectWithReasons|) devuelve también los
candidatos rechazados con su motivo de descarte. Los diez motivos posibles
se agrupan en \emph{estructurales} (incumplimiento de P1--P4 o ausencia de
patrón combinable) y de \emph{efecto colateral} (las cinco variantes de P5);
el catálogo completo está en el Apéndice~D y su árbol de decisión en la
\autoref{fig:arbol-p1-p5}. Como los motivos se recogen por candidato y se
agregan a nivel de método, un mismo método puede aportar varias categorías
a \lstinline|discardCategories|, lo que constituye evidencia agregada para
RQ1.

\subsection{Transformador}
La clase \texttt{NestedIfTransformer} implementa el algoritmo de
\autoref{sec:diseno}. Algunos detalles relevantes son: \textbf{clonación}
para no dejar el
AST inconsistente; \textbf{parentizado selectivo},
solo si la expresión es una expresión binaria con operador
\lstinline|OR|; \textbf{normalización del cuerpo}, produciendo siempre código con llaves;
\textbf{eliminación defensiva del \texttt{else}};
y \textbf{verificación previa}, replicando
P1 y P4 como salvaguarda.

\subsection{Cálculo proxy de complejidad cognitiva}
\label{sec:proxy}
La clase \texttt{CognitiveComplexityCalculator} calcula la complejidad
cognitiva de un método siguiendo el modelo de SonarSource
\cite{10.1145/3194164.3194186}. Implementa las tres reglas del modelo:

\begin{itemize}
  \item \textbf{Incremento estructural (+1)}: \lstinline|if|,
  \lstinline|else if|, \lstinline|else|, \lstinline|for|, \emph{for-each},
  \lstinline|while|, \lstinline|do-while|, \lstinline|switch|,
  \lstinline|catch|, \lstinline|break|/\lstinline|continue| con etiqueta y
  el operador ternario \lstinline|?:|.
  \item \textbf{Incremento de anidamiento (+nivel)}: \lstinline|if|,
  \lstinline|for|, \emph{for-each}, \lstinline|while|, \lstinline|do-while|,
  \lstinline|switch| y \lstinline|catch| (no en \lstinline|else if| ni
  \lstinline|else|). Además, los cuerpos de \emph{lambdas}, clases anónimas
  y clases locales se procesan a nivel de anidamiento $+1$.
  \item \textbf{Operadores lógicos}: $+1$ por cada secuencia de operadores
  del mismo tipo (\lstinline|&&| o \lstinline=||=) en \emph{cualquier}
  expresión (condiciones, \lstinline|return|, asignaciones, argumentos,
  etc.). El conteo recorre la expresión de forma plana ignorando por
  completo los paréntesis, tal como especifica el modelo de referencia, de
  modo que las secuencias alternantes (p.\,ej. \lstinline=a || b && c || d=)
  se contabilizan correctamente.
\end{itemize}

\paragraph{Limitaciones.}
La única regla del modelo no implementada es la \textbf{detección de recursión} (que
SonarSource penaliza con $+1$). Por lo demás, el cálculo cubre el modelo
completo.

\paragraph{Validación frente a los casos oficiales de SonarSource.}
La
corrección del proxy se comprueba a nivel unitario contra los casos de
prueba \emph{oficiales} del modelo: se ha implementado una clase que compara, con
aserciones de igualdad exacta, el valor calculado frente al valor de
referencia de SonarSource para cada caso del \emph{fixture}
\texttt{CognitiveComplexityMethodCheck}, el mismo conjunto que
emplea la herramienta \texttt{SoftwareCognitiveComplexityReducer}. Los
casos cubren operadores lógicos (incluidas las secuencias alternantes con
paréntesis ignorados), estructuras de control (\lstinline|switch|,
\lstinline|try|/\lstinline|catch|/\lstinline|finally|,
\lstinline|break|/\lstinline|continue| etiquetados,
\lstinline|if|/\lstinline|else if|/\lstinline|else| anidados y bucles) y
las características antes limitadas (ternario, \emph{lambda}, clase anónima
y clase local). El proxy reproduce el valor oficial en todos los casos
cubiertos; la única exclusión es la recursión.

\subsection{Pipeline experimental}
El pipeline experimental, implementado por \texttt{BatchRunner}, recorre el
corpus caso por caso y aplica las oportunidades \emph{de una en una} (la
política iterativa de la \autoref{sec:one-at-a-time}). Para no alterar el método original, mide su
complejidad cognitiva \emph{antes}, lo clona y transforma el clon, sobre el
que mide la complejidad \emph{después}; la diferencia entre ambas es el
$\Delta$ del caso. Además de transformar, anota \emph{por qué} no combinó
los condicionales que descartó: esos motivos de rechazo se recogen aunque
el caso acabe siendo elegible, porque constituyen la evidencia de RQ1. Como
cada pasada aplica una sola oportunidad y vuelve a detectar, el número de
oportunidades detectadas coincide con el de aplicadas, y se registra
cuántas pasadas hicieron falta (con un tope de seguridad de diez). Si un
fragmento no se puede analizar —porque no compila o no contiene ningún
método— se marca como no elegible, se le asigna una complejidad
\emph{antes} de $-1$ como valor centinela y se anota el motivo.

\subsection{Subsistema de evaluación de LLMs (RQ3)}
\label{sec:llm-subsistema}
El paquete \texttt{es.tfm.refactoring.llm} implementa el protocolo de RQ3.
La \autoref{fig:secuencia-rq3} resume el ciclo completo de una evaluación;
los apartados siguientes detallan cada pieza (prompt, \emph{provider},
oráculo y campaña).

\begin{figure}[H]
\centering
\includegraphics[width=0.92\linewidth]{fig-secuencia-rq3}
\caption[Secuencia de una evaluación RQ3]{Diagrama de secuencia de una
evaluación RQ3: del caso y su \emph{baseline} (de RQ2) al veredicto, pasando
por prompt, \emph{provider}, modelo y oráculo.}
\label{fig:secuencia-rq3}
\end{figure}

\paragraph{Prompt.} \texttt{LlmPromptBuilder} contiene literalmente la versión \lstinline|v1.0| del prompt: un \emph{system prompt} con las reglas de
combinación (P1--P4 reescritas como restricciones), exigencia de
preservación de orden de evaluación y semántica, prohibición de modificar
firmas o código fuera de los \lstinline|if|, y opción de rechazo explícito
mediante \lstinline|NO_REFACTORING_APPLICABLE|. El \emph{user prompt}
envuelve el código del caso en un bloque \lstinline|```java|.

\paragraph{Protocolo.} El protocolo (\texttt{LlmExperimentProtocol})
encapsula los parámetros como configuración inmutable. Declara los modelos
\lstinline|gpt-4o| y \lstinline|gpt-4.1|; temperatura \lstinline|0.0|; 3
intentos por caso; prompt \lstinline|v1.0|; y el máximo de tokens de salida
en 2048. Un protocolo auxiliar mantiene los parámetros pero
restringe la lista a \lstinline|gpt-4o|, y 
\emph{no sustituye} al protocolo formal de dos modelos. Existe además un
tercer protocolo, idéntico en parámetros
salvo la versión del prompt (\lstinline|v2.0|, \emph{few-shot}); donde su
finalidad es una comparación exploratoria \emph{zero-shot} (v1.0) frente a
\emph{few-shot} (v2.0) pero \textbf{no} forma parte de la campaña de
referencia.

\paragraph{Provider.} La interfaz \texttt{LlmResponseProvider} admite dos
implementaciones intercambiables: la invocación real a la API de OpenAI
(con credencial) y las respuestas grabadas para \emph{dry runs} y tests.

\paragraph{Oráculo}
\label{sec:oraculo}

La validación de las respuestas generadas por los modelos
(\texttt{LlmResponseValidator}) se implementa como una cadena de
comprobaciones secuenciales:

\begin{enumerate}
    \item \textbf{Detección de rechazo explícito.}  
    Se comprueba si la respuesta del modelo contiene el marcador \lstinline|NO_REFACTORING_APPLICABLE|.

    \item \textbf{Extracción del bloque de código.}  
    Se extrae el bloque \lstinline|```java| mediante una expresión regular.

    \item \textbf{Comprobación de parseabilidad.}  
    El código extraído se analiza con JavaParser para verificar que sea sintácticamente válido.

    \item \textbf{Existencia de al menos un método.}  
    Se comprueba que el fragmento parseado contenga al menos un método sobre el que aplicar la validación.

    \item \textbf{Preservación de la firma.}  
    Se verifica que la firma del método coincida con la del caso original. Si difiere, se registra como error de validación.

    \item \textbf{Medición de complejidad cognitiva \emph{after}.}  
    Una vez aceptado el código desde el punto de vista estructural, se calcula la complejidad cognitiva resultante.

    \item \textbf{Cálculo del delta y comparación con el baseline.}  
    Se obtiene el valor $\Delta$ y se compara con el delta del baseline determinista.

    \item \textbf{Asignación del veredicto final.}  
    A partir de las comprobaciones anteriores se asigna el veredicto correspondiente.
\end{enumerate}

El conjunto de veredictos (\texttt{LlmVerdict}) está formado por \texttt{SUCCESS}, \texttt{INCORRECT}, \texttt{INVALID\_OUTPUT}, \texttt{REFUSED}, \texttt{PARTIAL} y \texttt{ERROR}.

El oráculo incorpora además la \textbf{guarda de elegibilidad v1.1}: si el baseline declara que un caso no es elegible (porque no se cumple alguna de las precondiciones P1--P5) y el LLM aplica una transformación con \lstinline|delta != 0|, el veredicto pasa a ser \texttt{INCORRECT}, aunque la complejidad disminuya. Con esta comprobación se evita clasificar como aciertos transformaciones que violan las precondiciones del detector.

La \autoref{fig:oraculo} resume esta cadena de validación en forma de diagrama de flujo, mientras que la \autoref{tab:oraculo} recoge las reglas exactas de asignación de veredicto.

\begin{figure}[H]
\centering
\begin{tikzpicture}[
  term/.style={draw, rounded rectangle, fill=green!12, align=center,
               text width=3.1cm, inner sep=4pt, font=\small},
  endv/.style={draw, rounded rectangle, fill=gray!12, align=center,
               text width=3.3cm, inner sep=4pt, font=\small},
  proc/.style={draw, fill=blue!5, align=center, text width=4.8cm,
               inner sep=5pt, font=\small},
  dec/.style={draw, diamond, aspect=2.4, fill=orange!12, align=center,
              text width=2.5cm, inner sep=1pt, font=\scriptsize},
  arr/.style={-{Latex}}, lbl/.style={font=\scriptsize, inner sep=1.5pt},
]
\node[term] (start) at (0,0)     {Respuesta del LLM};
\node[dec]  (d1)    at (0,-2.3)  {\textquestiondown{}rechazo explícito?};
\node[proc] (p1)    at (0,-4.6)  {Extraer bloque \lstinline|```java|};
\node[dec]  (d2)    at (0,-6.9)  {\textquestiondown{}bloque, parsea y $\geq 1$ método?};
\node[proc] (p2)    at (0,-9.4)  {Comprobar firma; medir CC \emph{after}; $\delta$ y comparación con baseline};
\node[term] (end)   at (0,-11.7) {Asignar veredicto (Tabla~\ref{tab:oraculo})};
% Salidas laterales
\node[endv] (refv)  at (5.4,-2.3) {Veredicto de rechazo (Tabla~\ref{tab:oraculo})};
\node[endv] (inv)   at (5.4,-6.9) {\texttt{INVALID\_OUTPUT}};
% Flechas
\draw[arr] (start) -- (d1);
\draw[arr] (d1) -- node[lbl,left]{no} (p1);
\draw[arr] (p1) -- (d2);
\draw[arr] (d2) -- node[lbl,left]{sí} (p2);
\draw[arr] (p2) -- (end);
\draw[arr] (d1) -- node[lbl,above]{sí} (refv);
\draw[arr] (d2) -- node[lbl,above]{no} (inv);
\end{tikzpicture}
\caption{Cadena de validación del oráculo \texttt{LlmResponseValidator}
(diagrama de flujo). Las salidas tempranas cubren el rechazo explícito
(\lstinline|NO_REFACTORING_APPLICABLE|) y la salida no válida
(\texttt{INVALID\_OUTPUT}); el resto desemboca en la asignación de veredicto
según la \autoref{tab:oraculo}.}
\label{fig:oraculo}
\end{figure}

\begin{table}[H]
\centering
\caption{Reglas de asignación de veredicto del oráculo (política v1.1).
Fuente: \texttt{LlmResponseValidator.evaluate}. $\delta$ es la variación de
CC del LLM y \emph{baseline} la del prototipo determinista.}
\label{tab:oraculo}
\begin{tabular}{@{}p{0.46\linewidth}p{0.30\linewidth}l@{}}
\toprule
\textbf{Situación} & \textbf{Condición} & \textbf{Veredicto} \\
\midrule
Rechazo explícito, caso inelegible & rechazo $\wedge$ \lstinline|!eligible| & \texttt{SUCCESS} \\
Rechazo explícito, caso elegible & rechazo $\wedge$ \lstinline|eligible| & \texttt{REFUSED} \\
Sin bloque, no parsea o sin método & --- & \texttt{INVALID\_OUTPUT} \\
Firma modificada u otro error & errores $\neq \emptyset$ & \texttt{INCORRECT} \\
Inelegible y transforma (guarda v1.1) & \lstinline|!eligible| $\wedge\ \delta \neq 0$ & \texttt{INCORRECT} \\
Inelegible y no transforma & \lstinline|!eligible| $\wedge\ \delta = 0$ & \texttt{SUCCESS} \\
Elegible, $\delta<0$ igual o mayor que baseline & \lstinline|eligible| $\wedge\ \delta \le$ base $< 0$ & \texttt{SUCCESS} \\
Elegible, reduce menos que baseline & \lstinline|eligible| $\wedge$ base $< \delta < 0$ & \texttt{PARTIAL} \\
Elegible, no reduce ($\delta=0$, base $<0$) & \lstinline|eligible| $\wedge\ \delta = 0$ & \texttt{INCORRECT} \\
Error técnico (red, \emph{timeout}\dots) & excepción en el \emph{provider} & \texttt{ERROR} \\
\bottomrule
\end{tabular}
\end{table}

\paragraph{Campaña y exportación.} La campaña (\texttt{LlmCampaignRunner})
orquesta la ejecución sobre el subset, genera un registro por llamada
(timestamp, versión de modelo, prompt enviado, tiempo de respuesta) y delega
en el validador; las excepciones del \emph{provider} se propagan como
veredicto \texttt{ERROR}. El punto de
entrada ofrece \lstinline|--mode=live| (provider
real, validación de credenciales obligatoria) y \lstinline|--mode=dry-run|
(provider grabado, sin red).
